\section{Discussion}
\label{sec:discussion}

The cases support the paper's central claim in a specific form. Where a mechanical check
existed it worked, and where the plan was handed to a separate executor the handoff
exposed errors in the plan. What failed was what the original framework left implicit:
nobody recorded the reviewer's time, and nothing tied a reported number to the
computation that should have produced it. Both gaps concern the researcher's attention,
which is the quantity the framework puts at the center.

\subsection{Implications}
\label{sec:implications}

\paragraph{For researchers.} The assignment rule turns a vague question---how much should
I trust this agent?---into three questions about the task that can be answered before
any agent runs. Its monotonicity suggests the cheapest route to more delegation: change
the task rather than the model, by writing outputs to new files, adding a known-answer
test or keeping results local until they are accepted. Before a team is enlarged, the
review load $L$ and the number of defects found after acceptance $E$ should be measured on
a small pilot, because Equation~\eqref{eq:escape} shows that reading less is paid for in
escaped defects unless mechanical checks take up the difference. The source textbook
orders adoption accordingly: a group policy and a repository with an instruction file
come first, then one skill piloted at level 1 or 2, then checks and a reproducibility
package, and only then a small team, reviewed after 90 days on time saved, review time
and errors that escaped~\cite{mccreary2026aiforee}.

\paragraph{For advisors and research groups.} Most of the framework is a set of
conventions that a group can share: an instruction file, a vocabulary of statuses with
\textsf{accepted} reserved for people, a template for handoff records and a short record
of each acceptance before a result enters a paper. Case~B shows a side benefit: a status
field that agents cannot advance makes the review backlog visible, here as 50 simulations
awaiting approval. A backlog that is visible can be managed; one hidden in agents'
reports cannot. Because reviewers tend to accept automated output with less scrutiny
over time~\cite{parasuraman2010complacency,goddard2012automation}, a group can also ask
reviewers to state the expected value before reading an artifact or to check a random
sample rather than the first few items.

\paragraph{For builders of agent tools.} Four of the case findings point to features
of the harness rather than of the model. The harness, not the agent, should write
timestamps and record how long each artifact waited for and received review, so that
$k_j$, $r_j$ and $s_j$ can be measured. It should make \textsf{accepted} impossible for the
producing agent to set (I1). It should enforce step, time and spending limits and confirm
that a terminated process has stopped, which D11 shows an agent cannot be trusted to
report. And it should support the provenance rule by linking each number in a report to
the call that produced it, in the spirit of activity logging for
agents~\cite{chan2024visibility}.

\paragraph{For venues.} Current policies ask authors to disclose substantial use of
generative AI and to take responsibility for the
result~\cite{arxiv2025genai,nature2023tools,ec2025living}. A disclosure that named which
tasks were delegated, at which autonomy level, and which claims passed a human gate would
tell reviewers more than a statement that a tool was used. This is a suggestion that
follows from the framework, not a finding of this study.

\subsection{Limitations}
\label{sec:limitations}

\begin{enumerate}[leftmargin=*]
  \item \textit{Author bias.} The framework, the textbook it comes from and both real
    cases are the author's own, and one person coded the defects. An independent coder
    might classify some defects differently, particularly the boundary between agent
    checks and self-reports.
  \item \textit{Case selection.} Cases~B and~C are publication pipelines---a textbook and
    a paper---not laboratory research with instruments. The engineering case is
    illustrative. That the framework transfers to measurement-heavy research is argued
    from its structure, not shown.
  \item \textit{Unmeasured parameters.} Review times, the attention budget, the
    planning-to-execution ratio and the detection probabilities in
    Equations~\eqref{eq:load}--\eqref{eq:pipeline} are illustrative. The model's
    predictions, such as the span of control, have not been tested against observed
    teams.
  \item \textit{Binary task properties.} In practice reversibility, reach and
    verifiability are matters of degree: an action can be partly reversible, and a check
    can be cheap for one reviewer and expensive for another. The rule is a first cut that
    graded or probabilistic versions could refine.
  \item \textit{Incomplete records.} The logs were partly narrated by the agents
    themselves, timestamps were missing, and one log contained a number that no
    computation produced. Defects that were never recorded cannot be counted, so the 13
    defects are a lower bound, not a rate.
  \item \textit{One orchestrator.} The model assumes a single researcher. Groups with
    several reviewers, delegated review and shared queues raise coordination questions it
    does not address.
  \item \textit{Pace of change.} Agent products and models change within months. The
    framework is stated without reference to any product, but the parameter values that
    make it useful will drift and must be re-measured.
\end{enumerate}
